\documentclass[
    12pt,                % tamanho da fonte
    brazil               % idioma
    openany,             % capitulos podem comecar em qualquer pagina
    oneside,             % para impressao de um unico lado da folha
    chapter=TITLE,       % deixa a letra dos capitulos em maiuscula
    % sumario=tradicional, % estilo do sumario
    % fleqn
]{abntex2}
%
% ===============
%    PACKAGES    
% ===============
%
\usepackage[version=4]{mhchem}
\usepackage[scaled]{helvet}  % fonte nao serifada 
\usepackage{stix2}           % fonte serifada
\usepackage{graphicx}        % permite inserir imagens 
% \usepackage{microtype}      % melhora da tipografia (quebra o sumario)
\usepackage[T1]{fontenc}     % define melhorias para o tipo da fonte
\usepackage[utf8]{inputenc}  % codificacao do texto
\usepackage{indentfirst}     % identa o primeiro paragrafo
\usepackage{tabularray}      % permite criacao de tabelas personalizadas
\usepackage{xcolor}
\usepackage[locale=DE]{siunitx}
\usepackage[
    capitalise,              % nome das fig/tab/ em letras maiusculas
    nameinlink               % hyperref no nome das fig/tab
]{cleveref}
\usepackage{csquotes}
% nas versoes atuais do texlive, o sumario nao compila como deveria,
% entao a formatacao deve ser feita manualmente
% \usepackage{tocloft}
% \cftsetindents{part}{0em}{\cftlastnumwidth}
% \cftsetindents{chapter}{0em}{\cftlastnumwidth}
% \cftsetindents{section}{0em}{\cftlastnumwidth}
% \cftsetindents{subsection}{0em}{\cftlastnumwidth}
% \cftsetindents{subsubsection}{0em}{\cftlastnumwidth}
% \cftsetindents{paragraph}{0em}{\cftlastnumwidth}
% \cftsetindents{subparagraph}{0em}{\cftlastnumwidth}
% \renewcommand{\cftchapterfont}{\bfseries\MakeUppercase}
% \renewcommand{\cftchapterpagefont}{\normalsize\cftchapterfont}
% \renewcommand{\cftsectionfont}{\bfseries}
% \renewcommand{\cftsectionpagefont}{\cftsectionfont}
\renewcommand{\cftsubsectionfont}{\normalfont\sffamily}
\renewcommand{\cftsubsectionpagefont}{\cftsubsectionfont}
%
% ==================
%    REFERENCIAS    
% ==================
%
\usepackage[
    style=abnt,
    ittitles
]{biblatex}
\addbibresource{ref.bib}
%
% ==================
%    DOC OPTIONS    
% ==================
%
% as configuracoes a seguir formatam o estilo dos headings
\renewcommand{\ABNTEXchapterfont}{\sffamily\bfseries} % estilo do capitulo
\renewcommand{\ABNTEXchapterfontsize}{\Large}         % tamanho fonte capitulo
\setsecheadstyle{\sffamily\bfseries\large}            % estilo da secao
\setsubsecheadstyle{\sffamily\large}                  % estilo da subsecao
\setsubsubsecheadstyle{\sffamily\large}               % estilo da subsubsecao
\setlength{\afterchapskip}{\baselineskip}             % espacamento abaixo cap.
\UseTblrLibrary{booktabs}                             % para tabularray
\hypersetup{
    colorlinks=true,
    allcolors=blue!50!black,
    pdftitle={\@title},
    pdfauthor={\@author},
    pdfsubject={\imprimirpreambulo}
}
\captionnamefont{\footnotesize\sffamily\bfseries} % formatacao do titulo das figuras
\captiontitlefont{\footnotesize} % formatacao do texto do titulo das figuras
\newcommand{\ma}{\text{Ma}}
%
% ===========
%    CAPA    
% ===========
%
\renewcommand{\imprimircapa}{%
    \begin{center}\sffamily
        \includegraphics{./figuras/feis_unesp.pdf}
        \vfill
        {\bfseries\large AVALIAÇÃO DO MOTOR JOHN DEERE 6068 TIER 3}
        \vfill
        \begin{tabular}{ll}
            Nome & RA \\ \hline
            Alícia Ramos Modesto                  & 181052725 \\
            Gabriel Duarte da Silva               & 182054047 \\
            Higor Balsarini                       & 182055302 \\
            Matheus Henrique Panini               & 182053857 \\
            Yuri Fernando Oliveira Kazitani Cunha & 231052669
        \end{tabular}
        \vfill
        Docente: Gabriel Coelho Rodrigues Alvares
        \vfill
        Ilha Solteira -- SP \\
        Dezembro de 2023
    \end{center}
    \clearpage
}
%
\begin{document}
%
\imprimircapa\clearpage    % adicao da capa
\chapter*{}
\listoffigures\clearpage   % adicao da lista de figuras
\listoftables\clearpage    % adicao da lista de tabelas
\tableofcontents\clearpage % adicao do sumario

\textual

\chapter{Literature Review}\label{sec:literature_review}

This chapter deals with the history, the main concepts and some practical cases of \gls*{shm} inside the industry and academic area, besides showing how it may be used in the railway crack detection context.
Next, in the dynamic field, it will be studied the main mechanical concepts to get the necessary understanding to an \gls*{uav} motion as well the basics to know how an \gls*{uav} can be controlled.
Then, it will be shown the mathematics behind the algorithms of deep learning that will be implemented in the \cref{sec:results}. 
Finally, the way how the algorithms are going to be implemented and the tools necessary to achieve the desired \gls*{nn}.
\chapter{Metodologia}

\section{Revisão da Literatura}

Os conceitos teóricos explicados a seguir são baseados nas notas de aula, na obra de \textcite{heywood2018internal}, \textcite{ferguson2015internal}.

\subsection{Áreas}

Parâmetros geométricos são de suma importância para o motor. 
Áreas relacionadas às válvulas são necessárias para a compreensão do desempenho do motor.
Á área de cortina é a região ao redor da haste da válvula que é projetada para direcionar o fluxo de fluidos, como mostra a~\cref{fig:area_cortina}.
A área de garganta representa a seção mais estreita do canal de passagem, como mostra a~\cref{fig:area_garganta} influenciando a quantidade de massa fresca que entra na câmara de combustão.


\begin{figure}[!htb]
    \centering
     \begin{minipage}{0.49\textwidth}
        \centering
        \caption{Área de Cortina em uma Válvula}
        \includesvg[width=0.5\textwidth]{figuras/area_de_cortina.svg}
        \nota{A área hachurada representa a área de cortina.}
        \fonte{elaborado pelos autores.}
        \label{fig:area_cortina}
     \end{minipage}
     \hfill
     \begin{minipage}{0.49\textwidth}
        \centering
        \caption{Área de Garganta em uma Válvula}
        \includesvg[width=0.5\textwidth]{figuras/area_de_garganta.svg}
        \nota{A área hachurada representa a área de garganta.}
        \fonte{elaborado pelos autores.}
        \label{fig:area_garganta}
     \end{minipage}
   \end{figure}

\subsection{Coeficiente de Descarga}

O coeficiente de descarga é a razão entre o fluxo de ar que está passando através do componente durante o ensaio, pelo fluxo de ar que deveriapassar pelo componente durante o ensaio caso o escoamento isentrópico.
%
\begin{equation}
    C_d = \frac{V_r}{V_t}
    \label{eq:coef_descarga}
\end{equation}
%
onde $V_r$ o fluxo de massa real e $V_t$ é o fluxo de massa caso o escoamento fosse isentrópico.

\subsection{Pressão Média}

A pressão média indicada (IMEP) é aquela determinada baseada na geometria do motor e considerando as transformações do ciclo como irreversiveis. 
A pressão média de atrito (FMEP) indica a pressão perdida devido ao atrito e o bombeamento.
A pressão média efetiva (BMEP) é a pressão que o motor de fato tem trabalho líquido.
A BMEP pode ser calculada como mostra a seguir:
%
\begin{equation}
    \text{BMEP} = \frac{120 \dot{W}_b}{n V_d N}
    \label{eq:bmep}
\end{equation}
%
onde $\dot{W}_b$ e a potência, $n$ é quantidade de cilindros do motor, $V_d$ e o volume deslocado e $N$ é a rotação do motor em RPM.

\subsection{Carga Lateral}

A carga lateral do pistão é obtida através da análise de forças no pistão a partir das leis da mecânica dos sólidos e é dada por:
%
\begin{equation}
    F_L = \frac{\pi d P (R/L) \sin\alpha}{4\sqrt{1-(R/L)^2\sin^2\alpha}}
    \label{eq:carga_lateral}
\end{equation}
%
onde $d$ é o diâmetro do pistão, $P$ é a pressão exercida, $R$ é o comprimento da biela, $L$ é o comprimento da manivela e $\alpha$ é o ângulo percorrido pela manivela a partir do ponto morto superior (PMS).

\subsection{Velocidade do pistão}

A partir de geometria simples, é possível determinar a posição do cilindro em relação ao ângulo de rotação da manivela. Sabendo que $\alpha=\omega t$, determina-se a velocidade através da derivada da posicão do cilindro:
%
\begin{equation}
    v(t) = R \omega \sin{\left(\omega t \right)} + \frac{R^{2} \omega \sin{\left(\omega t \right)} \cos{\left(\omega t \right)}}{L \sqrt{1 - \dfrac{R^{2} \sin^{2}{\left(\omega t \right)}}{L^{2}}}}
    \label{eq:velocidade}
\end{equation}

\subsection{Número de Mach}

O número de Mach é um termo adimensional definido pela seguinte expressão:
%
\begin{equation}
    \ma = \frac{V}{c} = \frac{\text{velocidade do escoamento}}{\text{velocidade do som}}
    \label{eq:mach}
\end{equation}
%
que descreve a velocidade do escoamento. Quando $\ma = 1$ o escoamento é considerado sônico; quando $\ma < 1$, subsônico; quando $\ma > 1$, supersônico e quando $\ma \gg 1$, hipersônico \cite{cengel2015mecanica}.
\section{Montagem Experimental}

\section{Procedimento Experimental}
\chapter{Resultados e Discussão}

\section{Resultados medidos}

Para o ensaio de fluxo no cabeçote, foi necessário determinar os parâmetros geométricos relativos às válvulas, tanto de escape como a de admissão, visto que possuem diferentes dimensões; e coletar a vazão do fluxo de ar através da bancada de ensaios de fluxo.
Os parâmetros geométricos das válvulas de admissão e escape foram coletados e estão dispostos na~\cref{tab:dados_geometricos}. 
Os cálculos que envolvem esses parâmetros foram realizados utilizando a média aritmética de cada um deles.
Já com os ensaios realizados na bancada de ensaios de fluxo, foram obtidos os dados da~\cref{tab:dados_coletados_bancada}. Eles podem ser melhores visualizados na~\cref{fig:fluxo_levante}.
%
\begin{table}[!htb]
    \centering\footnotesize
    \caption{Parâmetros Geométricos das Válvulas}
    \begin{tblr}{
        row{odd} = azul!10, 
        row{1-2} = white,
        width=16cm,
        colspec = {*{6}{X[c]}}
    }
    \toprule
    \SetCell[c=3]{c} Admissão & & & \SetCell[c=3]{c} Exaustão & & & \\
    {$D_v$ \\ (mm)} & {$D_h$ \\ (mm)} & {$D_g$ \\ (mm)} & {$D_v$ \\ (mm)} & {$D_h$ \\ (mm)} & {$D_g$ \\ (mm)} \\
    \midrule
    24,88 & 4,93 & 20,48 & 28,44 & 5,10 & 24,60 \\
    25,10 & 4,90 & 19,85 & 28,40 & 5,00 & 24,35 \\
    24,90 & 4,90 & 20,50 & 28,45 & 4,95 & 24,40 \\
    24,47 & 5,00 & --- & 28,00 & 5,00 & --- \\
    \bottomrule
    \end{tblr}
    \nota{AVA: admissão na válvula de admissão, EVE: escape na válvula de escape, EVA: escape na válvula de admissão.}
    \fonte{elaborado pelos autores.}
    
    \label{tab:dados_geometricos}
\end{table}
%
\begin{table}[!htb]
    \centering\footnotesize
    \caption{Dados Coletados pela Bancada de Fluxo}
    \begin{tblr}{
        row{odd} = azul!10, 
        row{1-2} = {bg=white, font=\scriptsize},
        width=16cm,
        colspec = {*{9}{X[c]}}
    }
    \toprule
    \SetCell[c=3]{c} AVA & & & \SetCell[c=3]{c} EVE & & & \SetCell[c=3]{c} EVA & & & \\
     {Pressão \\ (in\ce{H2O})} & {Abertura \\ (mm)} & {Fluxo \\ (cfm)} & {Pressão \\ (in\ce{H2O})} & {Abertura \\ (mm)} & {Fluxo \\ (cfm)} & {Pressão \\ (in\ce{H2O})} & {Abertura \\ (mm)} & {Fluxo \\ (cfm)} \\
     \midrule
     25,00 & 0 & 0,00 & 25,00 & 0 & 0,00 & 25,00 & 0 & 0,00 \\
     25,03 & 1 & 29,11 & 25,01 & 1 & 25,97 & 25,01 & 1 & 24,12 \\
     25,03 & 2 & 54,12 & 25,02 & 2 & 53,92 & 25,01 & 2 & 49,43 \\
     25,02 & 3 & 74,55 & 25,02 & 3 & 75,77 & 25,02 & 3 & 72,81 \\
     25,04 & 4 & 87,17 & 25,01 & 4 & 91,36 & 25,02 & 4 & 87,03 \\
     25,03 & 5 & 94,37 & 25,02 & 5 & 105,32 & 25,02 & 5 & 92,68 \\
     25,04 & 6 & 98,87 & 25,04 & 6 & 115,26 & 25,02 & 6 & 94,28 \\
     25,04 & 7 & 101,04 & 25,03 & 7 & 119,16 & 25,02 & 7 & 95,16 \\
     25,02 & 8 & 102,35 & 25,02 & 8 & 120,74 & 25,02 & 8 & 95,86 \\
     24,99 & 9 & 103,70 & 24,92 & 9 & 121,74 & 25,02 & 9 & 96,33 \\
     25,01 & 10 & 104,07 & 25,04 & 10 & 121,68 & 25,01 & 10 & 96,74 \\
     \bottomrule
    \end{tblr}
    \nota{AVA: admissão na válvula de admissão, EVE: escape na válvula de escape, EVA: escape na válvula de admissão.}
    \fonte{elaborado pelos autores.}
    
    \label{tab:dados_coletados_bancada}
\end{table}
%
A partir da \cref{fig:fluxo_levante}, nota-se que, para pequenos levantes da válvula, o fluxo cresce proporcionalmente a abertura da válvula. 
Isto ocorre porque para pequenos levantes a área que limita o fluxo é a área de cortina. 
Contudo, ao abrir a válvula, a área que limitará o fluxo passará a ser a área de garganta.
Também pode ser visto que para a região de maiores levantes de válvula o fluxo tende a ser constante, mesmo com a contínua abertura da válvula.
Por outro lado, o gráfico também indica que, mesmo a válvula de admissão sendo maior que a de escape ela possui um fluxo menor. Isto ocorre porque na admissão os dutos são menores do que no escape, já que no escape necessita-se de maior facilidade para expelir os gases, enquanto na admissão deseja-se uma cinética nos gases de modo que, mesmo quando o pistão estiver no PMS, ainda seja possível colocar mais massa de ar no cilindro.
%
\begin{figure}[!htb]
    \centering
    \caption[Fluxo em Relação à Abertura da Válvula]{Fluxo em Relação à Abertura da Válvula. Os dados foram obtidos através do próprio software da bancada de fluxo Servitec WinSSFluxo.}
    \import{./figuras/}{levante_fluxo.pgf}
    \nota{AVA: admissão na válvula de admissão, EVE: escape na válvula de escape, EVA: escape na válvula de admissão.}
    \fonte{elaborado pelos autores.}
    \label{fig:fluxo_levante}
\end{figure}
%

Já em relação aos dados do motor John Deere, modelo 6068 Tier 3, os parâmetros geométricos foram retirados do manual do fabricante e os dados coletados foram retirados pelo sistema de controle durante o ensaio.

\section{Resultados calculados}

\subsection{Ensaio de Fluxo no Cabeçote}

\subsubsection*{Áreas}

A~\cref{fig:levante_ac} apresenta o comportamento da área da cortina em relação à abertura da válvula. Conforme ilustrado na~\cref{eq:area_cortina}, a área da cortina é diretamente proporcional ao levante da válvula. 
Dessa forma, à medida que a válvula aumenta sua abertura, a área da cortina também aumenta de maneira linear. 
A curva correspondente à válvula de admissão é mais íngreme devido ao maior diâmetro dessa válvula, que correspondete ao coeficiente angular da reta, como visto na~\cref{tab:dados_geometricos}.
%
\begin{figure}[!htb]
    \centering
    \caption[Área de Cortina em Relação à Abertura da Válvula]{Área de Cortina em Relação à Abertura da Válvula. A área de cortina aumenta gradativamente conforme a abertura da válvula aumenta, pois ela é uma função da abertura.}
    \import{./figuras/}{levante_ac.pgf}
    \nota{AVA: admissão na válvula de admissão, EVE: escape na válvula de escape, EVA: escape na válvula de admissão.}
    \fonte{elaborado pelos autores.}
    \label{fig:levante_ac}
\end{figure}
%
Em relação a área de cortina, percebe-se que ela cresce linearmente e tende a continuar a aumentar mesmo após ter superado a área de garganta. 
O fato de continuar abrindo a válvula mesmo após a área de garganta ser a limitante do fluxo, ocorre por alguns fatores, como por exemplo para evitar o fenômeno de flutuação da válvula. 
Para isso, ela continua abrindo, mas sua aceleração vai diminui, de modo que no seu retorno a pressão da mola consiga garantir que a válvula não flutue.

A área de garganta é fixa e foi calculada para as válvulas de admissão e escape, como mostra a~\cref{tab:area_garganta}.
%
\begin{table}[!htb]
    \centering\footnotesize
    \caption{Área de Garganta}
    \begin{tblr}{
        row{2} = azul!10,
        colspec={*{2}{X[c]}}
    }
    \toprule
    {Válvula de Admissão \\ (mm$^2$)} & {Válvula de Escape \\ (mm$^2$)} \\
    \midrule 
    610 & 465 \\
    \bottomrule
    \end{tblr}
    \fonte{elaborado pelos autores.}
    \label{tab:area_garganta}
\end{table}
%
\subsubsection*{Coeficiente de descarga}

O coeficiente de descarga leva em consideração a vazão real e a vazão isentrópica.
Para o cálculo da vazão teórica, diferentes abordagens podem ser realizadas.
Considerando a área de garganta e a área de cortina, no início da abertura da válvula, a área de cortina é muito menor que a área de garganta.
Conforme a abertura aumenta, a área de cortina aumenta, mas ela é fator limitante para o fluxo de ar.
Em um determinado instante, quando a área de garganta se iguala a área de cortina, então a área de garganta se torna um fator limitante, porque a área de cortina pode aumentar, mas a vazão não depende mais dela e sim da área de garganta.
A vazão, portanto, será limitada por uma das duas áreas, ou seja, no início do levante, a área de cortina será o fator limitante ($A_c < A_g$) até o momento em que as duas áreas se igualarem e então a partir daí a área de garganta será o fator limitante ($A_c > A_g$).
O coeficiente de descarga foi calculado a partir da~\cref{eq:coef_descarga_area_efetiva}, como visto na~\cref{fig:fluxo_cd}.
%
\begin{figure}[!htb]
    \centering
    \caption[Coeficiente de Descarga em Relação à Abertura da Válvula]{Coeficiente de Descarga em Relação à Abertura da Válvula. A figura a esquerda mostra $C_d$ em relação a área de garganta, enquanto a figura a direita mostra $C_d$ em relação a área mínima.}
    \import{./figuras/}{levante_cd_new.pgf}
    \nota{AVA: admissão na válvula de admissão, EVE: escape na válvula de escape, EVA: escape na válvula de admissão.}
    \fonte{elaborado pelos autores.}
    \label{fig:fluxo_cd}
\end{figure}
%
No primeiro gráfico estão  os valores do coeficiente de descarga considerando a área de garganta. 
Quando a válvula de admissão se abre, tem-se pequenos valores de coeficiente de descarga, pois a vazão é pequena quando comparada com a vazão da garganta para um escoamento isentrópico. 
Quando a abertura aumenta, o coeficiente de descarga aproxima-se da unidade, porém em nenhum momento ele a atinge. 
Isso ocorre, pois quando o levante de válvula aumenta, o escoamento se torna mais próximo ao escoamento da área de garganta, porém ele não chega à unidade, pois o escoamento que está ocorrendo na válvula não é um processo isentrópico, devido ao atrito e à irreversbilidade do sistema.
Já no segundo gráfico, coeficiente de descarga varia em relação a área miníma.
Ele representa variação real que acontece do coeficiente de descarga, visto que para pequenos levantes da válvula, área mínima é a área de cortina e, a medida que aumenta o levante de válvula, a área mínima passa a ser a área de garganta. 
Esta mudança de área mínima ocorre no levante de válvula de 5 mm para as condições de AVA e EVA, e 4 mm para EVE. Nestes pontos nota-se a mudança no gráfico, com as curvas tornando-se idênticas ao gráfico do lado esquerdo.

\subsection{Ensaio Dinamométrico}

A partir do ensaio realizado, sendo este monitorado e controlado pelo software, na sala de controle, uma planilha foi gerada com diversas informações em cada instante de tempo, para cada rotação em que o motor foi solicitado.
Considerando a variação da rotação, alguns parâmetros importantes foram coletados e dipostos na~\cref{fig:PTPotF_rot} em função da rotação, afim de observar o comportamento de cada um deles.

Enquanto a potência, torque e pressão foram coletados diretamente dos dados do software, a carga lateral foi calculada utilizando a~\cref{eq:carga_lateral},  considerando a pressão média efetiva, calculada através da~\cref{eq:bmep}.
Os pontos coletados são altamente variáveis, cada um com seu grau de variação, mesmo em um pequeno intervalo de variação da rotação. 
Portanto, foi feita a média dos valores coletados para cada intervalo de rotação e suas incertezas foram determinadas a partir do desvio padrão, portanto as curvas apresentam as barras de erro para melhor interpretação do resultados coletados.
%
\begin{figure}[!htb]
    \centering
    \caption{Principais Parâmetros em Função da Rotação}
    \import{./figuras/}{PTPotF_rot.pgf}
    \fonte{elaborado pelos autores.}
    \label{fig:PTPotF_rot}
\end{figure}
%
A manutenção do BMEP constante a partir de 1600 RPM, como mostra a \cref{fig:PTPotF_rot}, destaca a eficiência do motor na conversão de combustível em energia mecânica. 
Isso sugere um projeto otimizado, ajustado para condições operacionais específicas, maximizando a eficiência térmica e mecânica em uma faixa específica de rotações.
O torque constante, em torno de \SI{670}{N.m}, desempenha um papel importante na força de rotação do motor, impactando diretamente a aceleração e a resposta do veículo. Essa característica é útil em aplicações que demandam força constante, como em veículos pesados, proporcionando estabilidade em condições de carga variável.
A potência, como é esperado, varia quase que linearmente e aumenta gradativamente com o aumento da rotação.
A força lateral do pistão também contribui para a estabilidade global do motor, especialmente e estabiliza-se na faixa de 1600 RPM, assim como o BMEP, algo naturalmente esperado, visto que a força lateral é uma função dependente da BMEP.

\subsubsection*{Velocidade}

A velocidade média, como visto na~\cref{eq:velocidade_media}, varia linearmente com o curso do pistão, então naturalmente o gráfico da esquerda é uma reta que cresce linearmente a medida que a rotação também aumenta. 
%
\begin{figure}[!htb]
    \centering
    \caption{Velocidades em Função da Rotação}
    \import{./figuras/}{velocidades_rot.pgf}
    \fonte{elaborado pelos autores.}
    \label{fig:velocidades_rot}
\end{figure}
%
Já a velocidade instantânea, quando vizualizada em função do ângulo de manivela, tem o comportamento de uma senoide quando considerado seu comportamento padrão ideal, tendo picos de máximos e mínimos, que indicam a velocidade máxima do pistão em diferentes sentidos, isto é, quando sobe e quando desce.
No gráfico a direita, para um determinado valor de rotação, verifica-se quase uma linha vertical.
Representa o conjunto de velocidades instantâneas para aquela determinada rotação. 
Sendo assim, o valor mais alto e mais baixo em cada rotação representam as velocidades máximas, onde a abscissa representa a mudança de direção da velocidade.
Nota-se que com o aumento da rotação, a velocidade máxima do pistão auementa de forma gradativa e quase linear, atingindo \SI{15}{m.s^{-1}}, quando a rotação aproxima-se da 2000 RPM.

\subsubsection*{Pressão de Boost}

A pressão de boost é aquela relativa à pressão no coletor de admissão, e ocasionada pela presença do turbocompressor. 
Na~\cref{fig:pressao_boost_rot}, verifica-se o quanto a pressão de boost influencia principalmente na BMEP. 
No eixo esquerdo, com a curva em preto, está a pressão de boost; e no eixo da direita, com a curva em azul, está a BMEP. 
Duas ordenadas foram plotadas para melhor visualização da influência de uma na outra, pois a ordem de grandeza das duas são diferentes.
%
\begin{figure}[!htb]
    \centering
    \caption{Pressão de Boost em Função da Rotação}
    \import{./figuras/}{pressao_boost_rotacao.pgf}
    \fonte{elaborado pelos autores.}
    \label{fig:pressao_boost_rot}
\end{figure}
%
Observa-se que a pressão de boost corresponde a aproximadamente 10\% da BMEP e  ambas as curvas tem comportamento similar ao longo da variação da rotação. 
Este já é o comportamentamento esperado, visto que este motor é sobrecarregado, então a pressão de boost naturalmente fornece uma pressão maior e garante uma pressão média efetiva também maior.

\subsubsection*{Índice de Mach}

O índice de Mach considera a vazão de escoamento no motor em relação à velocidade do som.
Para a melhor eficiência volumétrica, deseja-se $\ma = 0,5$.
Duas abordagens foram realizadas, uma considerando a velocidade de escoamento isentrópica a partir da~\cref{eq:velocidade_isentropica} e aplicando na~\cref{eq:mach}; e outra utilizando os parâmetros geométricos do motor e a velocidade média, como visto na~\cref{eq:mach_f}. Os resultados estão dispostos na~\cref{tab:mach}.
%
\begin{table}[!htb]
    \centering
    \caption{Índice de Mach}
    \begin{tblr}{
        row{even} = azul!10,
        colspec={*{3}{X[c]}}
    }
    \toprule
     Método & Condição & Indice de Mach \\
     \midrule
     \cref{eq:mach} & Torque máximo & 0,59 $\pm$ 0,01 \\
     \cref{eq:mach} & Potência máxima & 2,7$\pm$ 0,2 \\
     \cref{eq:mach_f} & Torque máximo & 0,0808 $\pm$ 0,0001 \\
     \cref{eq:mach_f} & Potência máxima & 0,15022 $\pm$ 0,00015 \\
     \bottomrule
    \end{tblr}
    \fonte{elaborado pelos autores.}
    \label{tab:mach}
\end{table}
%
Considerando a~\cref{eq:mach}, verifica-se que para condições de potência máxima, a velocidade de escoamento não garante uma boa eficiência volumétrica, mas para condições de torque máximo, se aproxima do índice de Mach ideal.
Já pela~\cref{eq:mach_f}, para ambas as condições o índice de Mach se mostra bem baixo, revelando baixa eficiência volumétrica.

\chapter{Discussão e Conclusão}

\printbibliography
	
\end{document}